\documentclass[aspectratio=1610]{beamer}

\usepackage{amssymb}
\usepackage{amsmath}
\usepackage{amsfonts}
\usepackage{fitpptemplate}

\addbibresource{sources.bib}

\title{PA2: Jak své programování orientovat na objekty}
\author{Klub FIT++}
\date{2026}

\begin{document}

\maketitle

\section{PA2 naruby}

\begin{frame}{Co je jazyk C?}

\begin{itemize}
  \item Paradigmata
    \begin{itemize}
      \item Funkcionální (hlavním stavebním blokem C jsou funkce)
      \item Imperativní (popisujeme jak se co dělá, ne čeho chceme docílit)
      \item Procedurální (program rozdělujeme do procedur - funkcí - které se spouští v pořadí)\cite{geeks:procedural_language}
    \end{itemize}
  \item Relativně jednoduchý jazyk
    \begin{itemize}
      \item Hot take: vše je číslo, pole čísel, nebo n-tice čísel a polí
      \item Máme pár aritmetických / logických operací
      \item if/for/while na ovládání běhu
      \item Vše ostatní vychází z fungování OS/paměti (schováno ve funkcích)
    \end{itemize}
\end{itemize}

\end{frame}

\begin{frame}{Co tedy C++}

\begin{itemize}
  \item Prakticky rozšíření C (většinově funkční C je funkční C++)
  \item Drasticky komplikovanější (třídy, templates, exceptions, přetěžování, reference)
  \item C++ má vlastní ISO standard\cite{isocpp:n4860} - skoro 2000 stránek
  \item Paradigma Objektově Orientované Programování (OOP)
\end{itemize}

\end{frame}

\section{OOP}

\begin{frame}{Definujme si problém}

\begin{itemize}
  \item Řekněme, že tvoříme hru, potřebujeme uchovávat hráče
  \item Má počet životů (0 - 3), damage: 0 - inf, stav:
    \begin{itemize}
      \item Jumping
      \item Swimming
      \item OnGround
    \end{itemize}
\end{itemize}

\end{frame}

\begin{frame}[fragile]{Kód v C}

  \begin{columns}
    \begin{column}{0.35\textwidth}

      \begin{minted}{C}
      typedef enum {
        JUMPING,
        SWIMMING,
        ON_GROUND,
      } player_state_t;

      typedef struct {
        char *name;
        unsigned int health;
        unsigned int damage;
        player_state_t state;
      } player_t;
      \end{minted}
    \end{column}
    \begin{column}{0.75\textwidth}

        \pause
        \begin{minted}{C}
        int main() {
          player_t player;
          // What is player.health?...

          return 0;
        }
        \end{minted}

        \begin{minted}{C}
        int main() {
          player_t player = {0};
          // Player with 0 health??

          return 0;
        }
        \end{minted}

    \end{column}

  \end{columns}

\end{frame}

\begin{frame}[fragile]{Řešení v C}

\begin{minted}{C}

int create_player(player_t *player) {
  player.name = malloc(...);
  player.health = 3;
  player.damage = 5;
  player.state = ON_GROUND;

  return 0;
}

void delete_player(player_t *player) {
  free(player.name);
  player.name = NULL;
}

\end{minted}

\end{frame}


\begin{frame}[fragile]{Řešení v C}

\begin{minted}[fontsize=\small]{C}

int main(int argc, char **argv) {
  player_t player;
  if (!create_player(&player)) {
    fprintf(stderr, "Could not create player...");
    return -1;
  }

  // Kdokoliv může rozbít náš předpoklad
  player.health = UINT_MAX;

  // Zároveň musíme volat delete před každým return
  delete_player(&player);
  return 0;
}

\end{minted}

\end{frame}

\printbibliography

\end{document}
